% Part: normal-modal-logic
% Chapter: completeness
% Section: frame-completeness

\documentclass[../../../include/open-logic-section]{subfiles}

\begin{document}

\olfileid[id]{nml}{com}{fra}

\olsection{Kelengkapan Kerangka}

Teorema kelengkapan bagi \Log{K} dapat diperluas ke sistem modal lainnya
setelah kita menunjukkan bahwa model kanonik bagi suatu logika tertentu
mempunyai sifat kerangka yang bersesuaian.

\begin{thm}\ollabel{thm:completeframeprops}
  Jika suatu logika modal normal $\Sigma$ memuat salah satu !!{formula}s pada
  sisi kiri \olref{tab:correspondencetable}, maka model kanonik bagi~$\Sigma$
  mempunyai sifat yang bersesuaian pada sisi kanan.
\end{thm}

\begin{table}[htp]
  \centering
    \begin{tabular}{| l || l |}
      \hline
      \emph{Jika $\Sigma$ memuat \dots} & \emph{\dots model kanonik bagi
        $\Sigma$ bersifat:} \\
      \hline \hline
      \Ax{D}:\;\; $\Box!A \lif \Diamond !A$ & \quad serial; \\
      \hline
      \Ax{T}:\;\; $\Box!A \lif !A$ &\quad  refleksif;\\
      \hline
      \Ax{B}: \;\;$!A \lif \Box\Diamond!A$ &\quad  simetris; \\
      \hline
      \Ax{4}: \;\; $\Box!A \lif \Box\Box!A$ & \quad transitif; \\
      \hline
      \Ax{5}: \;\; $\Diamond !A \lif \Box\Diamond!A$& \quad Euklidean.\\
      \hline
    \end{tabular}
    \caption{Fakta korespondensi dasar.}
    \ollabel{tab:correspondencetable}
\end{table}

\begin{proof}
  Kita membahas masing-masing sifat secara berurutan.

  Andaikan $\Sigma$ memuat \Ax{D}, dan misalkan $\Delta \in W^\Sigma$; kita
  perlu menunjukkan bahwa terdapat suatu $\Delta'$ sedemikian sehingga
  $R^\Sigma \Delta\Delta'$. Cukup ditunjukkan bahwa $\Box^{-1}\Delta$
  $\Sigma$-konsisten, sebab berdasarkan Lema Lindenbaum kemudian terdapat
  suatu himpunan lengkap dan $\Sigma$-konsisten~$\Delta' \supseteq
  \Box^{-1}\Delta$, dan berdasarkan definisi $R^\Sigma$
  \iftag{prvBox}{}{serta \olref[mod]{lem:box-iff-diamond}, }berlaku
  $R^\Sigma \Delta\Delta'$. Jadi, demi memperoleh kontradiksi, andaikan
  $\Box^{-1}\Delta$ \emph{tidak} $\Sigma$-konsisten, yakni
  $\Box^{-1}\Delta \Proves[\Sigma] \lfalse$. Berdasarkan
  \olref[mod]{lem:box2}, $\Delta \Proves[\Sigma] \Box\lfalse$; dan karena
  $\Sigma$ memuat \Ax{D}, berlaku pula
  $\Delta \Proves[\Sigma] \Diamond\lfalse$. Namun, $\Sigma$ normal, sehingga
  $\Sigma \Proves \lnot\Diamond \lfalse$
  (\olref[prf][nor]{prop:notDiamondBot}); karena itu, berlaku pula $\Delta
  \Proves[\Sigma] \lnot\Diamond \lfalse$, bertentangan dengan konsistensi
  $\Delta$.

  Sekarang andaikan $\Sigma$ memuat \Ax{T}, dan misalkan $\Delta \in
  W^\Sigma$. Kita ingin menunjukkan bahwa $R^\Sigma \Delta\Delta$,
  yakni,\iftag{prvBox}{}{ berdasarkan \olref[mod]{lem:box-iff-diamond},}
  $\Box^{-1}\Delta \subseteq \Delta$. Namun, jika $\Box!A \in \Delta$,
  maka berdasarkan \Ax{T} berlaku pula $!A \in \Delta$, seperti yang
  diinginkan.

  Sekarang andaikan $\Sigma$ memuat \Ax{B}, dan andaikan $R^\Sigma
  \Delta\Delta'$ bagi $\Delta$, $\Delta' \in W^\Sigma$. Kita perlu
  menunjukkan bahwa $R^\Sigma \Delta'\Delta$, yakni,\iftag{prvBox}{
    $\Box^{-1}\Delta' \subseteq \Delta$. Berdasarkan
    \olref[mod]{lem:box-iff-diamond}, hal ini ekuivalen dengan}{}
  $\Diamond\Delta \subseteq \Delta'$. Jadi, andaikan $!A \in \Delta$.
  Berdasarkan \Ax{B}, berlaku pula $\Box\Diamond!A \in \Delta$. Berdasarkan
  hipotesis bahwa $R^\Sigma \Delta\Delta'$\iftag{prvBox}{}{ dan
    \olref[mod]{lem:box-iff-diamond}}, kita memperoleh
  $\Box^{-1}\Delta \subseteq \Delta'$, sehingga
  $\Diamond!A \in \Delta'$, seperti yang diperlukan.

  Sekarang andaikan $\Sigma$ memuat \Ax{4}, serta andaikan $R^\Sigma
  \Delta_1\Delta_2$ dan $R^\Sigma \Delta_2\Delta_3$. Kita perlu menunjukkan
  $R^\Sigma \Delta_1\Delta_3$. Dari hipotesis\iftag{prvBox}{}{ dan
    \olref[mod]{lem:box-iff-diamond}}, kita memperoleh sekaligus
  $\Box^{-1}\Delta_1 \subseteq \Delta_2$ dan
  $\Box^{-1}\Delta_2 \subseteq \Delta_3$. Untuk menunjukkan
  $R^\Sigma \Delta_1\Delta_3$, cukup ditunjukkan bahwa
  $\Box^{-1}\Delta_1 \subseteq \Delta_3$. Jadi, misalkan $!B \in
  \Box^{-1}\Delta_1$, yakni $\Box!B \in \Delta_1$. Berdasarkan \Ax{4},
  berlaku pula $\Box\Box!B \in \Delta_1$; dan berdasarkan hipotesis, pertama
  kita memperoleh $\Box!B \in \Delta_2$, lalu $!B \in \Delta_3$, seperti
  yang diinginkan.

  Sekarang andaikan $\Sigma$ memuat \Ax{5}, serta andaikan $R^\Sigma
  \Delta_1\Delta_2$ dan $R^\Sigma \Delta_1\Delta_3$. Kita perlu menunjukkan
  $R^\Sigma \Delta_2\Delta_3$. Hipotesis pertama menghasilkan
  $\Box^{-1}\Delta_1 \subseteq \Delta_2$\iftag{prvBox}{}{ berdasarkan
    \olref[mod]{lem:box-iff-diamond}}, sedangkan hipotesis kedua ekuivalen
  dengan $\Diamond\Delta_3 \subseteq \Delta_1$\iftag{prvBox}{,
    berdasarkan \olref[mod]{lem:box-iff-diamond}}{}. Untuk menunjukkan
  $R^\Sigma \Delta_2\Delta_3$\iftag{prvBox}{, berdasarkan
    \olref[mod]{lem:box-iff-diamond}, cukup}{, kita harus}
  menunjukkan bahwa $\Diamond\Delta_3 \subseteq \Delta_2$. Jadi, misalkan
  $\Diamond!A \in \Diamond\Delta_3$, yakni $!A \in \Delta_3$. Berdasarkan
  hipotesis kedua, $\Diamond!A \in \Delta_1$; dan berdasarkan \Ax{5},
  berlaku pula $\Box\Diamond!A \in \Delta_1$. Namun, kini hipotesis pertama
  menghasilkan $\Diamond!A \in \Delta_2$, seperti yang diinginkan.
  \end{proof}

Sebagai suatu akibat, kita memperoleh hasil kelengkapan bagi sejumlah sistem.
Sebagai contoh, kita mengetahui bahwa $\Log{S5} = \Log{KT5} =
\Log{KTB4}$ lengkap terhadap kelas semua model refleksif dan Euklidean, yang
sama dengan kelas semua model refleksif, simetris, dan transitif.

\begin{thm}\ollabel{thm:generaldet}
  Misalkan $\mClass{C}_\Ax{D}$, $\mClass{C}_\Ax{T}$,
  $\mClass{C}_\Ax{B}$, $\mClass{C}_\Ax{4}$, dan
  $\mClass{C}_\Ax{5}$ masing-masing merupakan kelas semua model serial,
  refleksif, simetris, transitif, dan Euklidean. Maka, bagi sebarang skema
  $!A_1$, \dots, $!A_n$ di antara \Ax{D}, \Ax{T}, \Ax{B}, \Ax{4}, dan
  \Ax{5}, sistem $\Log{K}!A_1 \dots !A_n$ ditentukan oleh kelas model
  $\mClass{C} = \mClass{C}_{!A_1} \cap \dots
  \cap \mClass{C}_{!A_n}$.
\end{thm}

\begin{prop}
  Misalkan $\Sigma$ suatu logika modal normal; maka:
  \begin{enumerate}
  \item\ollabel{prop:anotherfive-a}%
    Jika $\Sigma$ memuat skema $\Diamond!A \lif \Box
    !A$, maka model kanonik bagi $\Sigma$ bersifat fungsional parsial.
  \item Jika $\Sigma$ memuat skema $\Diamond!A \liff \Box
    !A$, maka model kanonik bagi $\Sigma$ bersifat fungsional.
  \item Jika $\Sigma$ memuat skema $\Box\Box!A \lif \Box
    !A$, maka model kanonik bagi $\Sigma$ bersifat rapat lemah.
  \end{enumerate}
(lihat \olref[frd][acc]{tab:anotherfive} untuk definisi sifat-sifat kerangka
  ini).
\end{prop}

\begin{proof}
  \begin{enumerate}
  \item Andaikan bahwa $\Sigma$ memuat skema $\Diamond !A \lif
    \Box !A$. Untuk menunjukkan bahwa $R^\Sigma$ bersifat fungsional parsial,
    kita perlu membuktikan bahwa bagi sebarang $\Delta_1$, $\Delta_2$,
    $\Delta_3 \in W^\Sigma$, jika $R^\Sigma \Delta_1\Delta_2$ dan $R^\Sigma
    \Delta_1\Delta_3$, maka $\Delta_2=\Delta_3$. Karena $R^\Sigma
    \Delta_1\Delta_2$, kita mempunyai $\Box^{-1}\Delta_1 \subseteq \Delta_2$;
    dan karena $R^\Sigma \Delta_1\Delta_3$, berlaku pula
    $\Box^{-1}\Delta_1 \subseteq \Delta_3$\iftag{prvBox}{}{, keduanya
      berdasarkan \olref[mod]{lem:box-iff-diamond}}. Identitas
    $\Delta_2=\Delta_3$ akan mengikuti jika kita dapat menetapkan kedua
    inklusi $\Delta_2 \subseteq \Delta_3$ dan $\Delta_3 \subseteq
    \Delta_2$. Bagi inklusi pertama, misalkan $!A \in \Delta_2$; maka
    $\Diamond!A \in \Delta_1$, dan berdasarkan skema serta ketertutupan
    deduktif $\Delta_1$, berlaku pula $\Box!A \in \Delta_1$. Karena hipotesis
    menyatakan bahwa $R^\Sigma \Delta_1\Delta_3$, diperoleh
    $!A \in \Delta_3$. Inklusi kedua serupa.

  \item Hal ini langsung mengikuti dari bagian~\olref{prop:anotherfive-a}
    dan pembuktian serialitas dalam \olref{thm:completeframeprops}.

  \item Andaikan $\Sigma$ memuat skema $\Box\Box!A \lif \Box!A$; untuk
    menunjukkan bahwa $R^\Sigma$ rapat lemah, misalkan $R^\Sigma
    \Delta_1\Delta_2$. Kita perlu menunjukkan bahwa terdapat suatu himpunan
    lengkap dan $\Sigma$-konsisten~$\Delta_3$ sedemikian sehingga
    $R^\Sigma \Delta_1\Delta_3$ dan $R^\Sigma \Delta_3\Delta_2$. Misalkan:
    \[
    \Gamma = \Box^{-1}\Delta_1 \cup \Diamond\Delta_2.
    \]
    Cukup ditunjukkan bahwa $\Gamma$ $\Sigma$-konsisten, sebab berdasarkan
    Lema Lindenbaum himpunan itu kemudian dapat diperluas menjadi suatu
    himpunan lengkap dan $\Sigma$-konsisten~$\Delta_3$ sedemikian sehingga
    $\Box^{-1}\Delta_1 \subseteq \Delta_3$ dan
    $\Diamond\Delta_2 \subseteq \Delta_3$, yakni
    $R^\Sigma \Delta_1\Delta_3$\iftag{prvBox}{}{ (berdasarkan
      \olref[mod]{lem:box-iff-diamond})} dan $R^\Sigma
    \Delta_3\Delta_2$\iftag{prvBox}{ (berdasarkan
      \olref[mod]{lem:box-iff-diamond})}{}.

    Demi memperoleh kontradiksi, andaikan $\Gamma$ tak konsisten. Maka
    terdapat !!{formula}s $\Box!A_1$, \dots, $\Box!A_n \in \Delta_1$
    dan $!B_1$, \dots,~$!B_m \in \Delta_2$ sedemikian sehingga
    \[!A_1, \dots, !A_n, \Diamond!B_1, \dots, \Diamond!B_m \Proves[\Sigma]
    \lfalse.\] Karena $\Diamond (!B_1 \land \dots \land !B_m) \to
    (\Diamond!B_1 \land \dots \land \Diamond!B_m)$ !!{derivable} dalam
    setiap logika modal normal, kita bernalar sebagai berikut dan memperoleh
    kontradiksi dengan konsistensi~$\Delta_2$:
    \begin{align*}
      !A_1, \dots, !A_n, & \Diamond!B_1, \dots,
      \Diamond!B_m \Proves[\Sigma] \lfalse \\
      !A_1, \dots,!A_n & \Proves[\Sigma] (\Diamond!B_1 \land
      \dots \land \Diamond!B_m) \lif \lfalse\\
      & \qquad\text{berdasarkan teorema deduksi}\\
      & \qquad\text{\olref[prf][prp]{prop:derivabilityfacts}\olref[prf][prp]{prop:derivabilityfacts-deduction}, dan \Taut} \\
      !A_1, \dots,!A_n & \Proves[\Sigma]
      \Diamond(!B_1 \land
      \dots \land !B_m) \lif \lfalse\\
      & \qquad \text{karena $\Sigma$ normal} \\
      !A_1, \dots,!A_n & \Proves[\Sigma]
      \lnot\Diamond (!B_1 \land \dots \land !B_m)\\
      & \qquad\text{berdasarkan \PL} \\
      !A_1, \dots,!A_n & \Proves[\Sigma]
      \Box\lnot (!B_1 \land \dots \land !B_m)\\
      & \qquad\text{$\Box\lnot$ untuk $\lnot\Diamond$} \\
      \Box!A_1, \dots,\Box!A_n
      & \Proves[\Sigma] \Box\Box \lnot (!B_1 \land \dots \land !B_m)\\
      & \qquad\text{berdasarkan \olref[mod]{lem:box1}} \\
      \Box!A_1, \dots,\Box!A_n
      & \Proves[\Sigma]
      \Box\lnot (!B_1 \land \dots \land !B_m)\\
      &\qquad\text{berdasarkan skema $\Box\Box!A \lif \Box!A$} \\
      \Delta_1 & \Proves[\Sigma]
      \Box\lnot (!B_1 \land \dots \land !B_m)\\
      &\qquad\text{berdasarkan monotonisitas, \olref[prf][prp]{prop:derivabilityfacts}\olref[prf][prp]{prop:derivabilityfacts-monotonicity}} \\
      & \Box\lnot (!B_1 \land \dots \land !B_m) \in \Delta_1\\
      &\qquad\text{berdasarkan ketertutupan deduktif}; \\
      & \lnot (!B_1 \land \dots \land !B_m) \in \Delta_2\\
      &\qquad \text{karena }
      R^\Sigma \Delta_1\Delta_2.
    \end{align*}
  \end{enumerate}
\end{proof}

Berdasarkan contoh-contoh ini, orang mungkin mengira bahwa setiap sistem
logika modal~$\Sigma$ bersifat \emph{lengkap}, dalam arti bahwa sistem itu
membuktikan setiap formula yang valid pada setiap kerangka tempat setiap
teorema dalam~$\Sigma$ valid. Sayangnya, terdapat banyak sistem yang tidak
lengkap dalam arti ini.

\end{document}
